% ===========================================================================
% RESULTS -- drop in after the setup section, before Discussion/Conclusion.
% Figures: fig4_models.pdf, fig5_tasks.pdf (in figures/).
% ===========================================================================

\section{Results}
\label{sec:results}
The fifteen complete configurations comprise 1,500 episodes, of which 285
(19.0\%) satisfy LIBERO's success predicate. The campaign consumed 83{,}153
model calls, 915{,}160 scored simulator steps, and 190.0 hours of API latency,
at a settled cost of \$3{,}655. Because every configuration covers the same ten
tasks $\times$ ten initial states, every comparison below is matched on task and
initial state by construction.

\subsection{Success by model and reasoning effort}
Figure~\ref{fig:models} gives the per-configuration success rate with its 95\%
Wilson interval.

\begin{figure}[t]
\centering
\includegraphics[width=\textwidth]{fig4_models.pdf}
\caption{Success rate for each model at each requested reasoning effort, over
100 episodes per configuration. Points are the observed rate and bars the 95\%
Wilson interval. Requested effort raises success for Astra, Sol, and Opus~5.5;
Opus~5 changes from 10\% to 12\% to 6\%. GPT-6 Luna records no success in 300
episodes.}
\label{fig:models}
\end{figure}

\placeholder{Three paragraphs to write here:}
\placeholder{(i) effort is monotone for Astra $0.28\to0.38\to0.44$, Sol
$0.24\to0.30\to0.33$, and Opus~5.5 $0.09\to0.18\to0.33$, while Opus~5 changes
$0.10\to0.12\to0.06$.}
\placeholder{(ii) accuracy and cost are close to independent --- Sol and
Opus~5.5 both reach 0.33 at high effort, at \$4.96 and \$5.46 per success
respectively, against \$94.57 for Opus~5 at high effort.}
\placeholder{(iii) Luna fails completely, and differently by effort: at low
effort it averages 8.5 calls and 137 steps and stops of its own accord, while
at medium and high it averages 67.5 and 74.0 calls and reaches the 80-call
allowance in 23\% and 45\% of episodes.}

\subsection{Success by task}
Figure~\ref{fig:tasks} pools the fifteen configurations to 150 episodes per
task.

\begin{figure}[t]
\centering
\includegraphics[width=\textwidth]{fig5_tasks.pdf}
\caption{Success rate by task, pooled over the fifteen complete
configurations. The three basket-packing tasks and the stove task carry nearly
all successes; three tasks are solved by no configuration at any effort
setting.}
\label{fig:tasks}
\end{figure}

\placeholder{One paragraph: the distribution is bimodal, and a task at 0/150
across five models and three effort levels makes an environment-side cause ---
an unreachable pose, an occluded target, or a predicate the achieved state does
not satisfy --- as plausible as a policy-side one. Inspect the rollouts for
tasks 4, 6 and 9 and say which the videos support.}

\subsection{Inference cost}
\placeholder{Table or short paragraph: calls per episode, tokens, latency, and
dollars per success per configuration. The headline is a factor-of-twenty
spread in cost per success across configurations of comparable accuracy.}

\subsection{How episodes end}
Of the 1,500 episodes, 285 (19.0\%) end in success, 1,115 (74.3\%) in an
unsuccessful terminal state, 93 (6.2\%) on a budget or limit stop, and 7
(0.5\%) on an API error. The limit stops are concentrated: 68 of the 93 belong
to GPT-6 Luna at medium and high effort. Episodes that stop on a limit or an
error are counted as failures in every rate above; none is silently dropped.

\placeholder{AUTHOR CHECK: the legacy \texttt{incomplete\_budget} label merges
the call limit and the dollar cap. Separate the 91 using the recorded call
count and settled cost --- ``the allowance bound the episode'' and ``spending
did'' have different implications.}
